\chapter{4x10\textsuperscript{9}D.C.}\label{4x10^9-7}

\hfill\break

Rimasi sul ponte mentre la \emph{Capetown Maru} lasciava l'attracco e si
dirigeva in mare aperto.

Non meno di una decina di navi portacontainer lasciarono Teluk Bayur,
mentre gli incendi divampavano, sgomitando per aggiudicarsi una
posizione migliore all'imbocco della baia. La maggior parte erano
piccole navi mercantili di dubbia provenienza, presumibilmente dirette a
Port Magellan nonostante ciò che dicevano i loro documenti di carico.
Erano imbarcazioni i cui proprietari e capitani avrebbero avuto molto da
perdere in caso di controlli.

Rimasi con Jala e ci reggemmo forte al corrimano, osservando un cargo
costiero pieno di ruggine che virava per uscire da un banco di fumo e si
stava avvicinando pericolosamente alla poppa della \emph{Capetown}.
Entrambi i natanti dettero l'allarme e l'equipaggio del ponte della
\emph{Capetown} guardò verso poppa con apprensione. Ma il cargo cambiò
direzione poco prima dell'impatto.

Quando fummo in alto mare, tra grosse onde di mareggiata e fuori dalla
protezione della baia, scesi sottocoperta per raggiungere Ina, Diane e
gli altri migranti in sala equipaggio. En sedeva a un tavolo montato su
cavalletti con i suoi genitori e Ibu Ina, e nessuno dei quattro sembrava
stare bene. Per rispetto delle condizioni di Diane, le era stata data
l'unica sedia imbottita della stanza, malgrado le sue ferite avessero
smesso di sanguinare e fosse riuscita a cambiarsi, indossando vestiti
asciutti.

Jala entrò nel salone un'ora dopo. Urlò per richiamare l'attenzione e
fece un discorso, che Ina tradusse per me: «Tralasciando il suo pomposo
autocompiacimento, Jala dice che è andato sul ponte di comando e ha
parlato col capitano. Tutti gli incendi scoppiati sul ponte sono spenti
e stiamo viaggiando in sicurezza. Il capitano si scusa per il mare
agitato. Secondo le previsioni meteo dovremmo superare la perturbazione
in tarda notte o domattina presto. Per le prossime ore, però\ldots»

Proprio in quel momento En, che era seduto accanto a Ina, si girò e le
vomitò addosso, in sostanza finendo la frase al posto suo.

ooo

Due notti dopo salii sul ponte con Diane a guardare le stelle.

Il ponte principale era più tranquillo di notte rispetto a qualunque ora
del giorno. Vi trovammo un posticino sicuro, tra i container da 40 piedi
e la sovrastruttura di poppa, dove potevamo parlare senza essere
ascoltati. Il mare era calmo, l'aria gradevolmente tiepida, e le stelle
brulicavano sulle ciminiere e i radar della \emph{Capetown} come fossero
impigliate nelle sartie.

«Stai ancora scrivendo il tuo memoriale?» Diane aveva visto
l'assortimento di memory card che avevo in valigia, oltre agli altri
supporti digitali e ai farmaci di contrabbando che avevamo portato da
Montreal. E poi taccuini, pagine staccate, appunti scarabocchiati.

«Non spesso come prima» risposi. «Non mi sembra più così urgente. Il
bisogno di appuntare tutto\ldots»

«Oppure la paura di dimenticare.»

«Oppure quello.»

«E ti \emph{senti} diverso?» chiese, sorridendo.

Ero un Quarto da poco tempo. Diane no. Oramai la sua ferita si era
rimarginata, senza lasciare altro che una striscia di carne raggrinzita
che seguiva la linea dell'anca. La capacità del suo corpo di rigenerarsi
da solo mi sorprendeva ancora una volta. Sebbene, con ogni probabilità,
l'avessi anch'io.

La sua domanda era un po' maliziosa. Avevo chiesto molte volte a Diane
se si sentisse diversa, da Quarto. La vera domanda era, ovviamente:
\emph{a me} pareva diversa?

Non c'era mai stata una risposta giusta. Chiaramente era una persona
diversa, dopo essere quasi morta e resuscitata nella Grande Casa - chi
non lo sarebbe stato? Aveva perso un marito e una fede e si era
risvegliata in un mondo che avrebbe fatto grattare la testa in segno di
perplessità persino a Buddha.

«La transizione è solo una porta» mi disse. «La porta d'ingresso di una
stanza. Una stanza dove non sei mai stato, anche se puoi averci dato una
sbirciata di tanto in tanto. Ora è la stanza in cui vivi; è la tua, ti
appartiene. Possiede caratteristiche che non puoi modificare - non puoi
renderla più grande o più piccola. Ma come la arredi dipende da te.»

«Sembra più un proverbio che una risposta» obiettai.

«Scusa. È il meglio che possa fare.» Si voltò verso le stelle. «Guarda,
Tyler, si vede l'Arco.»

Lo chiamavamo ``Arco'' per la tipica miopia dell'essere umano. In realtà
era un anello, un cerchio con un diametro di un migliaio e mezzo di
chilometri, di cui solo la metà si ergeva sul livello del mare. La parte
restante era sott'acqua oppure sprofondava nella crosta terrestre, e
forse (come alcuni avevano ipotizzato) sfruttava il magma suboceanico
come fonte energetica. Ma dal nostro punto di vista ad altezza di
formica era proprio un arco, la cui sommità si estendeva ben sopra
l'atmosfera.

Persino la metà emersa era visibile solo nelle fotografie scattate dallo
spazio, e anche quelle foto di solito erano manipolate per enfatizzare i
dettagli. Se avessimo potuto sezionare trasversalmente il materiale di
cui era composto l'anello - di fatto, il cavo che si piegava a cerchio -
avremmo visto un rettangolo di cinquecento metri nel lato più corto e di
un chilometro e mezzo in quello più lungo. Immenso, ma una frazione
infinitesimale dello spazio che racchiudeva e non sempre facile da
vedere da lontano.

La rotta della Capetown Maru ci aveva condotto a sud dell'anello,
paralleli al suo raggio e quasi in linea con l'apice. Il Sole splendeva
ancora sulla sommità, non era più una U o una J piegata ma una leggera
curvatura di ciglia (Diane la chiamava: ``Ciglio dello Stregatto'')
nell'alto del cielo del Nord. Le stelle gli roteavano alle spalle come
plancton fosforescente separato dalla prua di una nave.

Diane mi poggiò la testa sulla spalla. «Avrei voluto che Jason potesse
vedere tutto questo.»

«Credo che lo stia vedendo. Solo, non da questa prospettiva.»

ooo

Ci furono tre problemi immediati nella Grande Casa dopo la morte di
Jason.

Il più pressante fu Diane, le cui condizioni fisiche rimasero invariate
per giorni dopo l'iniezione del farmaco marziano. Era quasi in coma e
talvolta febbricitante, le pulsazioni sul collo come il battito d'ali di
un insetto. Eravamo a corto di materiale sanitario e dovetti convincerla
a bere un sorso d'acqua ogni tanto. L'unico vero miglioramento stava nel
suono del suo respiro, che si faceva sempre più rilassato e meno
apatico. Almeno, i polmoni stavano guarendo.

Il secondo problema fu sgradevole ma condiviso con troppe altre famiglie
del paese: un congiunto era morto e doveva essere sepolto.

Negli ultimi giorni, una grande ondata di morte (accidentale, suicida,
omicida) si era abbattuta sul mondo. Nessuna nazione sulla Terra era
preparata ad affrontarla, se non nel modo più spietato, e gli Stati
Uniti non facevano eccezione. Le radio locali iniziarono a comunicare
dove si trovavano i punti di raccolta delle fosse comuni; i camion
refrigerati erano stati sequestrati agli impianti di imballaggio della
carne; da quando il servizio telefonico era stato ripristinato, c'era un
numero da chiamare - ma Carol non voleva sentire ragioni. Quando
affrontai l'argomento, assunse una postura di fiera dignità: «Non lo
farò, Tyler. Non \emph{permetterò} che Jason venga gettato in un buco
come un poveraccio medievale.»

«Carol, non possiamo\ldots»

«Zitto» mi interruppe. «Ho ancora dei contatti risalenti ai vecchi
tempi. Fammi fare qualche telefonata.»

Un tempo era stata una stimata specialista e prima dello Spin doveva
aver avuto una grande rete di contatti; ma dopo trent'anni di isolamento
alcolico chi poteva conoscere? Ciononostante passò la mattinata al
telefono, rintracciando numeri che erano cambiati, ripresentandosi,
spiegando, persuadendo e supplicando. Tutto ciò mi sembrava un'impresa
disperata. Ma non più di sei ore dopo un carro funebre entrò nel
vialetto e due uomini, palesemente stremati ma sempre gentili e
professionali, entrarono, caricarono il corpo di Jason su una barella e
lo portarono fuori dalla Grande Casa per l'ultima volta.

Carol passò il resto della giornata al piano di sopra, tenendo la mano a
Diane e cantandole canzoni che probabilmente non poteva sentire. Quella
notte bevve il suo primo cocktail dalla mattina in cui il Sole rosso era
sorto - una ``dose di mantenimento'' la definì.

Il nostro terzo problema fu E.D. Lawton.

ooo

Bisognava mettere al corrente E.D. della morte del figlio e Carol si
preparò a occuparsi anche di quello. Confessò di non parlare con E.D.,
se non tramite avvocati, da un paio d'anni ormai, e che l'aveva sempre
terrorizzata, perlomeno da sobria. Lui era grosso, aggressivo e
minaccioso, mentre Carol era fragile, sfuggente e scaltra. Ma il suo
dolore aveva velatamente alterato l'equazione.

Ci vollero alcune ore ma alla fine riuscì a contattarlo - era a
Washington, ci avrebbe potuto raggiungere in breve tempo - e a dirgli di
Jason. Fu prudente e vaga sulle cause della morte. Gli spiegò che Jason
era arrivato a casa in preda a una polmonite o qualcosa di simile e e
che si era aggravato poco dopo che la corrente era saltata e il mondo
era impazzito - niente telefono, niente servizio ambulanze, in ultima
analisi nessuna speranza.

Le chiesi come aveva preso la notizia E.D.

Fece spallucce. «All'inizio non ha detto niente. Il silenzio è il modo
in cui E.D. esprime il dolore. Suo figlio è morto, Tyler. Può non averlo
sorpreso, considerato tutto quello che è successo negli ultimi giorni.
Ma gli ha fatto male. Credo che lo abbia ferito in maniera indicibile.»

«Gli hai detto che Diane è qui?»

«Ho pensato fosse più saggio non farlo.» Mi guardò. «Non gli ho detto
nemmeno che ci sei anche tu. So che Jason e E.D. erano ai ferri corti.
Jason era tornato a casa per scappare da qualcosa che stava succedendo
alla Perielio, qualcosa che lo spaventava, e che sospetto fosse in
qualche modo collegata al farmaco marziano. No, Tyler, non mi dare
spiegazioni - non voglio sentirle e probabilmente non le capirei. Ma ho
pensato che sarebbe stato meglio che E.D. non venisse a fare il
prepotente in giro per casa cercando di occuparsi di tutto.»

«Non ha chiesto di lei?»

«No, di Diane non ha parlato. C'è una cosa strana, però. Mi ha chiesto
di assicurarmi che Jason\ldots{} be', che il corpo di Jason fosse
salvaguardato. Mi ha tempestato di domande, a questo proposito. Gli ho
detto che mi sto organizzando, che ci sarà un funerale e che gli farò
sapere. Lui però non ha voluto lasciar cadere il discorso. Vuole
un'autopsia. Ma mi sono intestardita.» Mi guardò fredda. «Perché mai
vorrà un'autopsia, Tyler?»

«Non lo so» risposi.

Ma mi misi in testa di scoprirlo. Andai in camera di Jason, dove erano
state tolte le lenzuola dal letto ora vuoto. Aprii la finestra, mi
sedetti sulla sedia accanto alla credenza e guardai quello che si era
lasciato dietro.

Jason mi aveva chiesto di registrare le sue conoscenze conclusive sulla
natura degli Ipotetici e sulla loro manipolazione della Terra. Mi aveva
anche chiesto di inserire una copia di quella registrazione in una
decina di buste imbottite, munite di francobolli e indirizzi, e pronte
per essere spedite quando il servizio postale fosse ripreso. Chiaramente
Jase non si aspettava di produrre un tale monologo quando era arrivato
alla Grande Casa alcuni giorni dopo la fine dello Spin. Era perseguitato
da un altro tipo di emergenza. Il suo testamento sul letto di morte fu
un'appendice tardiva.

Sfogliai le buste. Gli indirizzi erano stati scritti a mano da Jason, e
appartenevano a nomi che non riconobbi. No, mi correggo; ne riconobbi
uno tra le buste.

Era il mio.

``\emph{Caro Tyler},

\emph{So che in passato ti ho oltremodo gravato di oneri. Mi dispiace ma
	sto per farlo di nuovo, e stavolta la posta in gioco è assai più alta.
	Lascia che ti spieghi. E scusami se ti sembrerò brusco, ma ho fretta,
	per motivi che ti appariranno chiari}.

\emph{I recenti episodi di ``sfarfallio'', come lo chiamano i media,
	hanno fatto scattare campanelli d'allarme nell'amministrazione Lomax. La
	stessa reazione che hanno provocato diversi altri avvenimenti, molto
	meno reclamizzati. Ti cito solo un esempio: dalla morte di Wun Ngo Wen,
	i campioni di tessuto presi dai suoi organi sono stati oggetto di studi
	presso il Centro Epidemiologico degli Animali a Plum Island, la stessa
	struttura dove era stato messo in quarantena quando era arrivato sulla
	Terra. La tecnologia marziana è subdola, ma la scienza forense moderna è
	cocciuta. Di recente è stato assodato che la fisiologia di Wun, in
	particolare il suo sistema nervoso, era stata alterata molto più di
	quanto la procedura della ``Quarta Età'' avesse illustrato nei suoi
	archivi. Per questa e altre ragioni, Lomax e la sua gente hanno
	cominciato a sentire puzza di bruciato. Hanno invitato E.D. a uscire dal
	suo riluttante pensionamento e stanno dando nuovo credito ai suoi
	sospetti sulle intenzioni di Wun. E.D. ha accolto con favore questa
	opportunità di rivendicare la Perielio (e la sua stessa reputazione), e
	non sta perdendo tempo per capitalizzare la paranoia della Casa Bianca}.

\emph{Come hanno deciso di procedere le autorità? Brutalmente. Lomax o i
	suoi consiglieri hanno ideato un piano per razziare le strutture
	esistenti della Perielio e confiscare tutto quello che abbiamo
	conservato dei beni e dei documenti di Wun, oltre alle nostre
	registrazioni e agli appunti di lavoro}.

\emph{E.D. non ha ancora collegato i puntini tra la mia guarigione dalla
	Sclerosi Multipla Atipica e i farmaci di Wun; o, se lo ha fatto, lo ha
	tenuto per sé. Oppure è ciò che preferisco credere. Perché se finisco
	nelle mani dei servizi di sicurezza la prima cosa che faranno sarà
	un'analisi del sangue e subito dopo diventerò un esperimento scientifico
	segregato, con ogni probabilità, nella vecchia cella di Wun a Plum
	Island. E non credo che E.D. lo voglia davvero. Per quanto possa
	portarmi rancore per avergli ``soffiato la Perielio'' o per aver
	collaborato con Wun Ngo Wen, rimane comunque mio padre}.

\emph{Ma non ti preoccupare. Anche se E.D. è tornato prepotentemente nel
	giro di Lomax alla Casa Bianca, ho le mie personali risorse. Le ho
	coltivate nel tempo. Queste persone non sono potenti nel senso comune
	del termine, sebbene alcune lo siano a modo loro, ma sono individui
	brillanti e rispettabili che hanno scelto di vedere il destino
	dell'umanità in una prospettiva a lungo termine, e grazie a loro sono
	stato avvertito in anticipo dell'irruzione alla Perielio. Ho attuato la
	mia fuga. Ora sono un latitante}.

\emph{Tu, Tyler, sei solo un presunto favoreggiatore, anche se alla fine
	potrebbe non cambiare la sostanza}.

\emph{Mi dispiace. Mi sento in parte responsabile di averti messo in
	questa posizione. Un giorno mi scuserò di persona. Per adesso posso solo
	offrirti un avvertimento}.

\emph{Le registrazioni digitali che ti ho consegnato quando hai lasciato
	la Perielio sono, ovvio, le parti censurate strettamente confidenziali
	prese dagli archivi di Wun Ngo Wen. Per quanto ne so potresti averle
	bruciate, sotterrate o gettate nell'oceano Pacifico. Non importa}.

\emph{Gli anni trascorsi a progettare navicelle spaziali mi hanno
	insegnato la virtù della ridondanza. Ho distribuito la saggezza di
	contrabbando di Wun a decine di persone in questo paese e in tutto il
	mondo. Non è ancora stata pubblicata in Internet - nessuno è tanto
	incosciente - ma è là fuori. Questo è senza dubbio un atto profondamente
	antipatriottico e criminale. E se sarò catturato verrò accusato di
	tradimento. Nel frattempo sto cercando di trarne il massimo vantaggio}.

\emph{Ma non credo che una simile conoscenza (che comprende protocolli
	per la modificazione umana che, tra le altre cose, possono curare gravi
	patologie e di cui so qualcosa) debba essere chiusa in un cassetto per
	interessi nazionali, anche se la sua divulgazione pone ulteriori
	problemi}.

\emph{Lomax e il suo Congresso addomesticato sono ovviamente contrari.
	Perciò ho intenzione di disperdere gli ultimi frammenti dell'archivio e
	tagliare la corda. Mi darò alla macchia. Dovresti farlo anche tu. In
	effetti, dovrai farlo. Chiunque abbia avuto a che fare con me, incluso
	tutto il personale della vecchia Perielio, prima o poi è destinato a
	finire sotto controllo federale}.

\emph{Oppure, al contrario, potrai volerti presentare al più vicino
	ufficio dell'FBI e consegnare il contenuto di questa busta. Se è ciò che
	credi più opportuno segui la tua coscienza; non ti biasimerò, ma non
	posso garantire per ciò che potrebbe scaturirne. La mia esperienza con
	l'amministrazione Lomax mi suggerisce che la verità non ti renderà
	libero, anzi}.

\emph{In ogni caso mi rincresce metterti in una posizione difficile. Non
	è giusto. È chiedere troppo a un amico, e sono sempre stato orgoglioso
	di chiamarti amico}.

\emph{Forse E.D. aveva ragione su una cosa. La nostra generazione ha
	lottato per trent'anni per recuperare ciò che lo Spin ci aveva tolto
	quella notte di ottobre. Ma non possiamo. Non c'è niente in questo
	universo in evoluzione a cui aggrapparsi, e non si ottiene niente a
	provarci. Se ho imparato qualcosa dalla mia ``Quartità'', è proprio
	questo. Siamo effimeri come gocce di pioggia. Tutti cadiamo, e finiamo a
	terra da qualche parte}.

\emph{Cadi libero, Tyler. Usa i documenti allegati se ti servono. Mi
	sono costati cari ma sono totalmente affidabili}.

(\emph{È bello avere degli amici ai piani alti}!)''

I ``documenti allegati'' erano, in sostanza, un pacchetto di documenti
d'identità di riserva: passaporti, carte d'identità della Sicurezza
Nazionale, patenti di guida, certificati di nascita, numeri di
previdenza sociale, persino diplomi di laurea in medicina. Tutti
contenevano la mia descrizione, ma nessuno aveva il mio vero nome.

ooo

La fase di ripresa di Diane stava procedendo. Il polso si era rafforzato
e i polmoni ripuliti, sebbene avesse ancora la febbre. Il farmaco
marziano stava facendo il suo dovere, la ricostruiva dall'interno,
rielaborava e modificava il suo DNA nei modi più ingegnosi.

Mentre le sue condizioni miglioravano, cominciò a fare domande prudenti
- sul Sole, sul Pastore Dan, sul viaggio dall'Arizona alla Grande Casa.
A causa della febbre intermittente, le risposte che le davo non le
rimanevano impresse. Mi chiese più di una volta cosa fosse successo a
Simon. Se era lucida le raccontavo del vitello rosso e del ritorno delle
stelle; se era stordita le dicevo solo che Simon era da un'altra parte e
che sarei tornato a prendermi cura di lei più tardi. Nessuna di queste
risposte - le verità o le mezze verità - sembravano soddisfarla.

Alcuni giorni era apatica, stava appoggiata alla finestra e guardava la
luce del Sole che scorreva lungo le lenzuola avvallate. Altri giorni era
febbrilmente irrequieta. Un pomeriggio chiese un foglio e una
penna\ldots{} ma quando glieli diedi tutto ciò che scrisse fu una
singola frase: ``\emph{Sono forse io il custode di mio fratello}'',
ripetuta fino a farsi venire i crampi alle dita.

«Le ho detto di Jason» ammise Carol quando le mostrai il foglio.

«Sei sicura che sia stata una scelta saggia?»

«Doveva saperlo prima o poi. Se ne farà una ragione, Tyler. Non ti
preoccupare. Diane starà bene. Diane è sempre stata quella forte.»

ooo

La mattina del giorno del funerale preparai le buste che Jason aveva
lasciato, introducendo in ognuna una copia della sua ultima
registrazione, attaccai i francobolli e le infilai in una cassetta della
posta scelta a caso sulla via che portava alla cappella locale dove
Carol aveva fissato il servizio funebre. I pacchetti avrebbero dovuto
attendere qualche giorno prima di venire raccolti - il servizio postale
era ancora in via di ripristino - ma immaginai che fossero più al sicuro
lì che nella Grande Casa.

La ``cappella'' era un'impresa di pompe funebri ecumenica su una strada
principale suburbana, densa di traffico ora che le restrizioni sui
viaggi erano state revocate. Jase aveva sempre provato uno sdegno
razionalista nei confronti dei funerali elaborati, ma il senso di
dignità di Carol richiedeva una cerimonia, per quanto insulsa e
\emph{pro forma}. Era riuscita a radunare una piccola folla, in gran
parte composta da vicini di lunga data che si ricordavano di Jason da
piccolo e avevano intravisto la sua carriera nelle brevi dichiarazioni
televisive e nelle colonne laterali dei quotidiani. Era il suo status di
celebrità che svaniva a riempire le panche.

Pronunciai un breve elogio funebre. (Diane avrebbe fatto di meglio, ma
era troppo malata per partecipare.) «Jase,» dissi, «ha dedicato tutta la
vita alla ricerca della conoscenza, senza arroganza e con umiltà: aveva
capito che la conoscenza non è stata creata ma scoperta; non poteva
essere posseduta, solo condivisa, passata di mano in mano, di
generazione in generazione. Jason si è reso parte integrante di tale
condivisione e ne fa ancora parte. Ha intessuto se stesso nella tela del
sapere.»

E.D. entrò nella cappella mentre ero ancora sul pulpito.

Era a metà della navata quando mi riconobbe. Mi fissò per un lungo
minuto prima di accomodarsi nella panca libera più vicina.

Era più smunto di quanto mi ricordassi, e aveva rasato gli ultimi
capelli bianchi riducendoli a un'invisibile peluria. Ma aveva ancora il
portamento di un uomo di potere. Indossava in modo impeccabile un
completo su misura. Incrociò le braccia e ispezionò la stanza con aria
imperiale, controllando chi era presente. Il suo sguardo si soffermò su
Carol.

Quando la funzione finì, Carol rimase con coraggio a ricevere le
condoglianze che le fecero i vicini mentre uscivano. Aveva pianto molto
negli ultimi giorni ma in quel momento sembrava determinata a non farlo,
distaccata in modo quasi clinico. E.D. le si avvicinò dopo che fu uscito
l'ultimo ospite. Lei si irrigidì come un gatto che percepisce la
presenza di un predatore più grande.

«Carol,» esordì E.D., «Tyler.»

Mi lanciò uno sguardo aspro.

«Nostro figlio è morto» disse Carol. «Jason se n'è andato.»

«Per questo sono qui.»

«Spero che tu sia qui per piangerlo\ldots»

«Certo.»

«\ldots{} e non per qualche altro motivo. Perché lui è tornato a casa
per fuggire da te. Suppongo che tu lo sappia.»

«So molto più di quanto immagini. Jason era confuso\ldots»

«Era molte cose, E.D., ma non era confuso. Ero insieme a lui quando è
morto.»

«C'eri davvero? Interessante. Perché, a differenza tua, io ero insieme a
lui quand'era vivo.»

Carol tirò un forte sospiro e voltò la testa di scatto come se fosse
stata schiaffeggiata.

E.D. disse: «Andiamo, Carol. Sono stato io a crescerlo, e lo sai bene.
Potrai non essere d'accordo con il tipo di vita a cui l'ho destinato, ma
è ciò che ho fatto - dargli una vita e i mezzi per viverla.»

«Io l'ho dato alla luce.»

«Quella è una funzione fisiologica, non un atto morale. Tutto ciò che
abbia mai posseduto l'ha avuto da me. Tutto quello che ha imparato
gliel'ho insegnato io.»

«Nel bene e nel male\ldots»

«E ora mi vuoi biasimare solo perché ho delle preoccupazioni
pratiche\ldots»

«\emph{Quali} preoccupazioni pratiche?»

«Ovviamente sto parlando dell'autopsia.»

«Sì. Me l'hai fatto presente al telefono. Ma è indecoroso e francamente
impossibile.»

«Speravo che avresti preso sul serio le mie preoccupazioni. Chiaramente
non lo hai fatto. Ma non mi serve il tuo permesso. Fuori da questo
edificio ci sono degli uomini che aspettano di reclamare il corpo, e ti
possono mostrare un mandato ai sensi della Legge sulle Misure
d'Emergenza.»

Fece un passo indietro allontanandosi da lui. «Hai così tanto potere?»

«Né tu né io abbiamo voce in capitolo. Succederà, che ci piaccia o no.
Ed è davvero una pura formalità. Non gli sarà fatto alcun male. Quindi,
per l'amor di Dio, non perdiamo dignità e rispetto reciproco. Lasciami
prendere il corpo di mio figlio.»

«Non posso farlo.»

«Carol\ldots»

«Non posso darti il suo corpo.»

«Non mi stai ascoltando. \emph{Non hai scelta}.»

«No, scusami, sei \emph{tu} che non stai ascoltando \emph{me}. Senti,
E.D. \emph{Non} posso darti il suo corpo.»

Lui aprì la bocca e poi la richiuse. Gli occhi si spalancarono.

«Carol,» sibilò, «cos'hai fatto?»

«Non c'è \emph{nessun} corpo. Non più.» Le labbra di lei si curvarono in
un sorriso scaltro e amaro. «Ma suppongo che tu possa prendere le sue
ceneri. Se insisti.»

ooo

Accompagnai Carol alla Grande Casa, dove il suo vicino, Emil Hardy - che
non appena era stata ripristinata la corrente aveva smesso di pubblicare
il giornaletto locale dalla vita breve - era rimasto seduto con Diane.

«Abbiamo parlato dei vecchi tempi nel quartiere» disse Hardy mentre
andava via. «Ricordo quando eravate bambini e vi guardavo correre sulle
bici. Era tanto tempo fa. Questo \emph{problema} che ha\ldots»

«Non è contagioso» lo rassicurò Carol. «Non ti preoccupare.»

«È strano, però.»

«Sì. È strano davvero. Grazie Emil.»

«A me e Ashley farebbe piacere se veniste a cena qualche volta.»

«Mi pare una bella cosa. Ringrazia Ashley da parte mia per favore.»
Chiuse la porta e si rivolse a me: «Mi serve qualcosa da bere. Ma
andiamo per gradi. E.D. sa che sei qui. Perciò te ne devi andare e devi
portare Diane via con te. Puoi farlo? Portarla in un posto sicuro? Dove
E.D. non possa trovarla?»

«Certo che posso. Ma che ne sarà di te?»

«Non sono in pericolo. E.D. potrebbe mandare qui della gente a cercare
un qualche tesoro che Jason, secondo lui, potrebbe avergli rubato. Ma
non troverà nulla - ammesso che tu sia \emph{scrupoloso}, Tyler - e non
può togliermi la casa. Io e E.D. abbiamo firmato un armistizio molto
tempo fa. Le nostre schermaglie sono irrilevanti. Ma può fare del male
\emph{a te e} può farlo a Diane, anche se non ne ha intenzione.»

«Non permetterò che accada.»

«Allora raccogli le tue cose. Potreste non avere molto tempo.»

ooo

Il giorno prima che la \emph{Capetown Maru} attraversasse l'Arco secondo
l'agenda stabilita, salii sul ponte per guardare il sorgere del Sole. La
volta era per lo più invisibile, i suoi pilastri sprofondati erano
nascosti dall'orizzonte a est e a ovest, ma nella mezz'ora prima
dell'alba l'apice rimase allineato nel cielo quasi sopra le nostre
teste, affilato come una lama e delicatamente brillante.

A metà mattinata era svanito dietro al velo di un alto cirro, ma noi
sapevamo che c'era.

La prospettiva del transito rendeva tutti nervosi - non solo i
passeggeri ma anche l'equipaggio esperto. Tutti continuarono a svolgere
le loro mansioni abituali, prendendosi cura della nave, aggiustando i
macchinari, scalpellando e ridipingendo le strutture di sostegno, ma nel
loro ritmo di lavoro c'era una bruschezza che il giorno prima non avevo
notato. Jala salì sul ponte portandosi dietro una sedia di plastica e
venne dov'ero seduto io, protetto dal vento dai container da 40 piedi ma
con una vista ristretta del mare.

«Questo è il mio ultimo viaggio dall'altra parte» mi confidò.

Per il tepore del giorno, indossava dei jeans e una camicia svolazzante
gialla. L'aveva sbottonata per esporre il petto alla luce del Sole.
Prese una lattina di birra dal distributore e strappò la linguetta.
Tutte queste azioni dicevano di lui che era un uomo di mondo, un uomo
d'affari, che disprezzava allo stesso modo la \emph{sharia} musulmana e
l'\emph{adat} dei Minangkabau. «Stavolta,» disse, «non si torna
indietro.»

Aveva bruciato i ponti dietro di sé - letteralmente, se aveva avuto a
che fare con la rivolta orchestrata a Teluk Bayur. (Le esplosioni
avevano garantito una copertura sospettosamente opportuna per la nostra
fuga, sebbene fossimo stati quasi coinvolti nella deflagrazione.) Per
anni Jala aveva portato avanti un commercio di intermediazione per il
traffico degli immigrati molto più redditizio della sua legittima
azienda di import-export. Si facevano più soldi con la gente che con
l'olio di palma, diceva. Ma la concorrenza tra India e Vietnam era dura
e il clima politico si era inasprito; a quel punto era meglio ritirarsi
a Port Magellan che passare il resto dei propri giorni in una galera
della nuova Reformasi.

«Hai mai fatto il transito?»

«Due volte.»

«È stato difficile?»

Scrollò le spalle. «Non credere a tutto ciò che senti.»

Prima di mezzogiorno molti passeggeri erano sul ponte. Oltre ai
Minangkabau del villaggio, c'erano vari migranti acehnesi, malesi e
tailandesi, in tutto un centinaio - troppo numerosi per le cabine a
disposizione, per questo tre container di alluminio, accuratamente
ventilati, erano stati attrezzati come dormitori.

Questa non era la truce e spesso mortale tratta degli esseri umani, che
di solito trasportava i rifugiati in Europa o in Nord America. La
maggior parte delle persone che attraversavano l'Arco ogni giorno erano
esuberi delle fragili Nazioni Unite - inclusi in programmi di
reinsediamento approvati e sovvenzionati. Eravamo trattati con rispetto
dall'equipaggio, molti di loro avevano trascorso mesi a Port Magellan e
ne avevano capito lusinghe e tranelli.

Un mozzo aveva fatto spazio in una zona del ponte principale
trasformandola in una sorta di campo di calcio, limitato da alcune reti,
in cui giocava un gruppo di bambini. Ogni tanto la palla rimbalzava
fuori, spesso addosso a Jala, con suo grande disappunto. Quel giorno era
irritabile.

Gli chiesi quando la nave avrebbe fatto il transito.

«Secondo il capitano, a meno che non cambiamo velocità, dodici ore o giù
di lì.»

«Il nostro ultimo giorno sulla Terra» dissi.

«Non scherzare.»

«Intendevo in senso letterale.»

«E abbassa la voce. I marinai sono superstiziosi.»

«Che cosa farai a Port Magellan?»

Jala inarcò le sopracciglia. «Che cosa farò? Scoperò donne bellissime. E
qualcuna brutta, di sicuro. Che altro?»

Il pallone da calcio rimbalzò di nuovo oltre le reti. Stavolta Jala lo
raccolse e se lo strinse contro la pancia. «Maledizione, vi avevo
avvertito! Questa partita è finita!»

Una decina di bambini si lanciarono subito contro le reti, protestando a
tutta voce, ma fu En che raccolse il coraggio di avvicinarsi per
affrontare Jala di petto.

Il ragazzo sudava, la cassa toracica pompava come un mantice. La sua
squadra era in vantaggio di cinque goal. «Ridaccelo, per favore» disse.

«Lo vuoi indietro?» Jala si alzò in piedi stringendo forte il pallone,
arrogante e misteriosamente arrabbiato. «Lo vuoi? Vallo a prendere.»
Dette un calcio alla palla che le fece descrivere una lunga traiettoria
e la portò oltre il ponte, dentro l'immensità verde acqua dell'oceano
Indiano.

En sembrò prima sbalordito poi furioso. Borbottò qualcosa di volgare e
aspro in minang.

Jala arrossì. Poi gli tirò un ceffone a mano aperta, talmente forte da
fargli saltare gli occhiali pesanti, che scivolarono sul pavimento del
ponte.

«Chiedi scusa» pretese Jala.

En cadde su un ginocchio, gli occhi serrati. Fece un paio di respiri
singhiozzanti. Alla fine si alzò. Avanzò di qualche passo sulle tavole
del ponte e recuperò gli occhiali. Li rimise a posto e tornò indietro
con quella che mi parve un'incredibile dignità. Si parò di fronte a
Jala.

«No» reagì con un filo di voce. «\emph{Tu} chiedi scusa.»

Jala sussultò e imprecò. En si fece piccolo piccolo. Jala alzò di nuovo
la mano.

Gli bloccai il polso a mezz'aria.

Jala mi guardò, sbigottito. «Cosa fai! Lasciami.»

Provò a liberare la mano. Non glielo permisi. «Non lo colpire ancora» lo
ammonii.

«Faccio quello che mi pare!»

«Va bene,» dissi, «ma non lo colpire di nuovo.»

«\emph{Tu}\ldots{} dopo quello che ho fatto per te!»

Poi mi lanciò un'altra occhiataccia.

Non so cosa vide sul mio viso. Non so con esattezza cosa stessi provando
in quel momento. Qualunque cosa fosse, sembrò confonderlo. Il suo pugno
chiuso si allentò. Parve appassire.

«Maledetto pazzo americano» borbottò. «Vado alla mensa.» E poi aggiunse,
rivolto alla piccola folla di bambini e marinai che si erano radunati
attorno a noi: «Dove c'è \emph{pace} e \emph{rispetto}!» E se ne andò.

En mi stava ancora fissando, a bocca aperta.

«Mi dispiace per ciò che è successo» gli dissi.

Annuì.

«Non posso recuperare il tuo pallone» aggiunsi.

Si toccò la guancia dove Jala lo aveva schiaffeggiato. «Non c'è
problema» rispose debolmente.

Più tardi - durante la cena alla mensa dell'equipaggio, a qualche ora
dall'attraversamento - raccontai l'incidente a Diane. «Non ho pensato a
quello che facevo. Mi pareva\ldots{} scontato. Quasi un riflesso. È una
cosa da Quarti?»

«Potrebbe esserlo. L'istinto di proteggere una vittima, specialmente un
bambino, e farlo in un attimo, senza pensare. L'ho provato anch'io.
Presumo sia qualcosa che i marziani hanno scritto nella loro
riorganizzazione neurale\ldots{} supponendo che possano davvero
progettare sentimenti così impercettibili. Vorrei che Wun Ngo Wen fosse
qui per spiegarcelo. Oppure Jason, del resto. Ti è sembrata una
forzatura?»

«No\ldots»

«O sbagliato, inappropriato?»

«No\ldots{} credo che fosse proprio la cosa giusta da fare.»

«Ma prima del trattamento non lo avresti fatto?»

«Forse sì. O avrei voluto. Ma probabilmente sarei rimasto nel dubbio
finché non fosse stato troppo tardi.»

«Quindi non sei dispiaciuto per il tuo comportamento.»

No. Solo sorpreso.

Era dipeso tanto da me quanto dalla biotecnologia marziana, questo mi
stava dicendo Diane, e supposi che fosse vero\ldots{} ma ci sarebbe
voluto un po' per abituarmici. Come ogni altro tipo di transizione
(dall'infanzia all'adolescenza, dall'adolescenza all'età adulta) c'erano
nuovi imperativi con cui avere a che fare, nuove opportunità e insidie,
nuovi dubbi.

Per la prima volta nel giro di molti anni ero di nuovo uno sconosciuto
per me stesso.

ooo

Avevo quasi finito di preparare i bagagli quando Carol scese al piano di
sotto, leggermente ubriaca, traballante, con una scatola da scarpe in
mano. La scatola aveva l'etichetta {RICORDI} ({SCUOLA}).

«Devi prendere questa» si raccomandò. «Era di tua madre.»

«Se per te ha un significato profondo, Carol, tienila.»

«Grazie, ma ho già preso ciò che volevo.»

Aprii il coperchio e guardai il contenuto. «Le lettere.»

Le lettere anonime indirizzate a Belinda Sutton, il nome da nubile di
mia madre.

«Sì. Quindi le avevi viste. Le hai mai lette?»

«No, non proprio. Quanto bastava per capire che erano lettere d'amore.»

«Oddio. Suona così sdolcinato. Preferisco pensarle come tributi. Sono
molto caste, davvero, se le leggi con attenzione. Non sono firmate. Tua
madre le ha ricevute quando eravamo entrambe all'università. All'epoca
frequentava tuo padre e non poteva certo mostrargliele - anche lui le
scriveva delle lettere. Per questo le condivise con me.»

«Non ha mai scoperto chi le aveva scritte?»

«No. Mai.»

«Doveva essere curiosa.»

«Certo. Ma era già fidanzata con Marcus a quel tempo. Cominciò a uscire
con Marcus Dupree quando lui e E.D. stavano lanciando la loro prima
azienda, progettazione e manifattura di palloni da alta quota, ai tempi
in cui gli aerostati erano tecnologia ``blue sky'', come la chiamava
Marcus: un po' folli, un po' idealisti. Belinda aveva soprannominato
Marcus e E.D. ``i fratelli Zeppelin''. Quindi credo che noi fossimo le
sorelle Zeppelin, Belinda e io. Perché lì cominciai a flirtare con E.D.
In un certo senso, Tyler, il mio matrimonio non è stato altro che un
tentativo di continuare ad avere tua madre come amica.»

«Le lettere\ldots»

«Interessante, no, che le abbia conservate per tutti questi anni? Alla
fine le chiesi perché. Perché non buttarle via? Mi rispose: ``Perché
sono sincere.'' Fu il suo modo di onorare chiunque le avesse scritte.
L'ultima arrivò una settimana prima del suo matrimonio. Poi
nessun'altra. E un anno dopo sposai E.D. Anche come coppie eravamo
inseparabili, non te l'ha mai detto? Facevamo le vacanze insieme,
andavamo al cinema assieme. Belinda venne in ospedale quando nacquero i
gemelli e io aspettavo sulla porta quando \emph{ti} portò a casa per la
prima volta. Ma tutto questo finì quando Marcus ebbe l'incidente. Tuo
padre era un uomo meraviglioso, Tyler, molto concreto, molto divertente
- l'unica persona che riusciva a far ridere E.D. Spericolato
all'inverosimile, però. Belinda fu devastata dalla sua morte. E non solo
emotivamente. Marcus aveva sperperato quasi tutti i loro risparmi e
Belinda spese ciò che rimaneva per gestire l'ipoteca sulla loro casa a
Pasadena. Perciò quando E.D. si trasferì a Est e facemmo un'offerta per
questa casa ci venne naturale invitarla a usare la casa degli ospiti.»

«In cambio delle pulizie» aggiunsi.

«Quella fu un'idea di E.D. Io volevo solo averla vicino. Il mio
matrimonio non era riuscito come il suo. Al contrario. In quel periodo
Belinda era la mia unica amica, più o meno. Quasi una confidente.» Carol
sorrise. «Quasi.»

«È per questo che vuoi tenere le lettere? Perché fanno parte della tua
storia con lei?»

Mi sorrise come si fa a un bambino ritardato. «No, Tyler. Non hai
capito. \emph{Sono mie}.» Il suo sorriso si assottigliò. «Non stupirti.
Tua madre era tanto eterosessuale e senza complicazioni quanto qualsiasi
altra donna che abbia mai conosciuto. Ho solo avuto la sfortuna di
innamorarmi di lei. Me ne innamorai in modo così abietto che avrei fatto
qualunque cosa - persino sposare un uomo che sembrava, già dal
principio, un po' sgradevole - solo per tenermela vicino. E durante
tutto quel tempo, Tyler, in tutti quegli anni di silenzio non le ho mai
detto ciò che provavo. Mai, tranne che in quelle lettere. Ero contenta
che le avesse tenute, anche se mi sembravano un po' pericolose, come
qualcosa di esplosivo o radioattivo nascosto in bella vista. Erano la
prova della mia stessa stupidità. Quando tua madre morì - intendo il
giorno stesso della sua morte - andai nel panico; provai a nascondere la
scatola; pensai di distruggere le lettere ma non ne fui capace, non
riuscii a costringermi a farlo; e poi, dopo il divorzio con E.D., quando
non dovevo più ingannare nessuno, le tenni semplicemente per me stessa.
Perché, vedi, sono \emph{mie}. Sono sempre state mie.»

Non sapevo cosa dire. Carol vide la mia espressione e scosse la testa
con tristezza. Mi mise le fragili mani sulle spalle. «Non essere
arrabbiato. Il mondo è pieno di sorprese. Siamo nati tutti estranei a
noi stessi e agli altri, e raramente veniamo presentati ufficialmente.»

ooo

Quindi passai quattro settimane nella stanza di un motel nel Vermont a
prendermi cura di Diane durante la sua guarigione.

La sua guarigione fisica, per meglio dire. Il trauma emotivo subito al
ranch Condon e in seguito l'aveva stremata e fatta isolare. Diane aveva
chiuso gli occhi su un mondo che sembrava alla fine e li aveva riaperti
quando la realtà aveva ormai perso i suoi punti cardinali. Non era in
mio potere sistemare le cose.

Perciò cercai di rendermi utile ma con cautela. Spiegai ciò che doveva
essere spiegato. Non feci alcuna richiesta e chiarii che non mi
aspettavo alcuna ricompensa.

Il suo interesse verso il mondo cambiato si risvegliò per gradi. Chiese
del Sole, restituito al suo aspetto benevolo, e le riportai ciò che mi
aveva detto Jason: la membrana dello Spin era ancora al suo posto; anche
se la camera temporale era finita, continuava a proteggere la Terra così
come aveva sempre fatto, modificando le radiazioni letali in un
simulacro di luce solare accettabile per l'ecosistema del pianeta.

«Allora perché l'hanno spento per sette giorni?»

«Lo hanno \emph{abbassato}, non \emph{spento} del tutto. E l'hanno fatto
per consentire che qualcosa attraversasse la membrana.»

«Quella cosa nell'oceano Indiano.»

«Sì.»

Mi chiese di farle ascoltare le registrazioni delle ultime ore di Jason,
e nel sentirle pianse. Mi chiese delle sue ceneri. Le aveva portate via
E.D. o le aveva tenute Carol? (Nessuna delle due. Carol mi aveva
consegnato l'urna e mi aveva detto di disporne in qualunque maniera
ritenessi appropriata. «L'orribile verità, Tyler, è che tu lo conoscevi
meglio di me. Jason per me era un codice cifrato, il figlio di suo
padre. Invece, per te era un amico.»)

Osservammo il mondo che riscopriva se stesso. Le sepolture di massa alla
fine cessarono; i sopravvissuti in lutto e spaventati iniziarono a
rendersi conto che il pianeta aveva riacquistato un futuro, per quanto
strano sembrasse. Per la nostra generazione fu una stupefacente
inversione. Il manto dell'estinzione ci era caduto dalle spalle; cosa
avremmo fatto senza di esso? Che cosa avremmo fatto, ora che non eravamo
più spacciati ma semplicemente mortali?

Vedemmo i filmati, provenienti dall'oceano Indiano, della mostruosa
struttura che si era conficcata nella carne del pianeta, l'acqua marina
che evaporava nei punti di contatto con gli enormi pilastri. L'Arco,
come lo cominciarono a chiamare tutti, abbreviando il nome Portale ad
Arco, nato dalla sua forma e dalle navi che avevano fatto ritorno con
storie di boe di navigazione scomparse, condizioni meteo inusuali,
bussole che giravano impazzite e una linea costiera selvaggia dove non
avrebbe dovuto esserci alcun continente. Varie flotte vennero
prontamente inviate. Il testamento di Jason aveva accennato a una
spiegazione, ma solo in pochi avevano avuto il vantaggio di sentirla -
io, Diane e la decina di persone che l'avevano ricevuta per posta.

Cominciò a fare esercizi un po' alla volta, ogni giorno, andava a
correre su un sentiero sterrato dietro al motel quando faceva più
fresco, tornando con l'odore di foglie morte e di legna bruciata tra i
capelli. Il suo appetito migliorò. E anche il menu nella caffetteria: la
consegna dei beni alimentari era stata ripristinata; l'economia
nazionale si stava rimettendo in moto scricchiolando.

Apprendemmo che anche Marte era stato dissigillato. Erano passati dei
segnali tra i due pianeti; il Presidente Lomax, in uno dei suoi discorsi
del tipo ``tutti radunati attorno alla bandiera'', aveva persino
lasciato intendere che il programma spaziale con equipaggio sarebbe
stato ripristinato, un primo passo verso l'instaurazione di rapporti
costanti con quello che chiamava (con sospetta esuberanza) ``il nostro
pianeta affine''.

Parlammo del passato. Parlammo del futuro.

Ci conoscevamo troppo bene, o non abbastanza. Avevamo un passato ma
nessun presente. E Diane era divorata dall'ansia per la sparizione di
Simon fuori Manassas.

«Ti ha quasi lasciato morire» le ricordai.

«Non di proposito. Non è malvagio. Tu lo sai.»

«Allora è pericolosamente ingenuo.»

Diane chiuse gli occhi pensierosa. Poi disse: «C'è una frase che il
Pastore Bob Kobel amava usare quando eravamo al Jordan Tabernacle: ``Il
suo cuore gridava a Dio''. Se quella frase descrive qualcuno, quello è
Simon. Ma io ci ho riflettuto. Quel cuore che grida credo sia qualcosa
di universale, che si riferisca a tutti noi. Tu, Simon, io, Jason.
Persino Carol. Persino E.D. Quando ci rendiamo conto di quanto sia
grande l'universo e breve una vita umana, i nostri cuori gridano. A
volte è un grido di gioia: credo fosse così per Jason; penso sia questo
che non ho capito di lui. Aveva il dono dello stupore. Ma per quasi
tutti noi è un grido di terrore. Il terrore dell'estinzione o
dell'insignificanza. I nostri cuori gridano. Forse a Dio, o forse solo
per rompere il silenzio.» Si scostò i capelli dalla fronte e vidi che il
braccio, prima pericolosamente magro, era di nuovo rotondo e forte.

«Penso che il grido uscito dal cuore di Simon fosse il suono umano più
puro del mondo. Vedi, lui non è bravo a capire le persone; è ingenuo; ed
è il motivo per cui ha girato così tanti tipi di fede, il Nuovo Regno,
il Jordan Tabernacle, il ranch Condon\ldots{} qualunque congrega,
bastava che parlasse chiaro e rispondesse al bisogno di dare un senso
all'esistenza.»

«A costo della tua vita?»

«Non ho detto che è saggio. Sto dicendo che non è cattivo.»

In seguito sono arrivato a riconoscere quel tipo di discorso: parlava
come un Quarto. Distaccato ma impegnato. Intimo ma oggettivo. Non mi
dispiaceva, ma a volte mi faceva drizzare i peli sul collo.

ooo

Non molto dopo averla dichiarata completamente guarita, Diane mi informò
che voleva andarsene. Le chiesi dove volesse andare.

Doveva trovare Simon, disse. Doveva ``sistemare le cose'', in un modo o
nell'altro. Dopotutto erano ancora sposati. Era importante, per lei,
sapere se fosse vivo o morto.

Le ricordai che non aveva soldi da spendere né un posto suo dove stare.
Disse che se la sarebbe cavata in qualche modo. Allora le diedi una
delle carte di credito che mi aveva fornito Jason, avvisandola che non
potevo garantire per la carta - non avevo idea di chi pagasse il premio,
quale fosse il limite di credito o se qualcuno potesse risalire lei,
tracciandola.

Mi chiese come avrebbe potuto rimanere in contatto con me.

«Chiamami e basta» risposi. Aveva il mio numero, quello che avevo pagato
e mantenuto invariato per molti anni, collegato a un telefono che mi
portavo sempre dietro malgrado non suonasse quasi mai.

Poi la accompagnai alla stazione locale degli autobus, dove sparì in
mezzo a una folla di turisti dispersi, rimasti bloccati dalla fine dello
Spin.

ooo

Il telefono suonò sei mesi dopo, quando i giornali stavano ancora
pubblicando grandi titoli sul ``nuovo mondo'' e i canali delle
televisioni via cavo avevano cominciato a mostrare filmati di un
promontorio roccioso e selvaggio ``da qualche parte oltre l'Arco''.

Fino ad allora centinaia di imbarcazioni grandi e piccole avevano fatto
la traversata. Alcune erano grandi spedizioni scientifiche, approvate
dall'IGY o dall'ONU, scortate da navi americane e corredate di agenzie
di stampa. Altre erano barche private. Altri ancora erano pescherecci a
strascico che tornavano in porto con il loro carico di pescato, simile
al merluzzo se visto sotto una luce soffusa. Tutto questo era
severamente vietato, certo, ma prima che venisse dichiarato l'embargo il
``merluzzo dell'Arco'' si era infiltrato in tutti i maggiori mercati
asiatici. Si rivelò commestibile e nutriente. Cosa che, come avrebbe
potuto dire Jase, era un indizio: quando il pesce fu sottoposto
all'analisi del DNA il genoma suggerì una remota discendenza terrestre.
Il nuovo mondo non era soltanto ospitale, sembrava essere stato
allestito pensando all'umanità.

«Ho trovato Simon» mi informò Diane.

«E quindi?»

«Vive in un parcheggio per camper fuori Wilmington. Racimola un po' di
soldi facendo riparazioni domestiche - biciclette, tostapane, quel
genere di cose. Peraltro prende un sussidio e frequenta una piccola
chiesa pentecostale.»

«È stato felice di vederti?»

«Non riusciva a smettere di scusarsi per ciò che era successo al ranch
Condon. Diceva che avrebbe voluto farsi perdonare. Mi ha chiesto se
poteva fare qualcosa per semplificarmi la vita.»

Strinsi un po' più forte il telefono. «Cosa gli hai risposto?»

«Che voglio il divorzio. Era d'accordo. E ha detto altro. Ha detto che
ero cambiata, che c'era qualcosa di diverso in me. Non riusciva a capire
cosa. Ma non credo che abbia gradito.»

``L'odore di zolfo, forse.''

«Tyler,» mi domandò Diane, «sono cambiata così tanto?»

«Tutto cambia» le risposi.

ooo

L'altra chiamata importante da parte sua arrivò un anno dopo. Ero a
Montreal, grazie anche alla carta d'identità contraffatta di Jason,
aspettavo che mi venisse riconosciuto lo status di immigrato e lavoravo
in una clinica ambulatoriale a Outremont.

Dall'ultima conversazione con Diane, le dinamiche di base dell'Arco
erano state comprese. I fatti erano contraddittori per chiunque
concepisse l'Arco come una macchina statica o una semplice ``porta'', ma
guardandola come faceva Jason - una complessa e consapevole entità in
grado di percepire e manipolare gli eventi all'interno del suo dominio -
aveva molto più senso.

Due mondi erano stati collegati tramite l'Arco, ma solo per le
imbarcazioni oceaniche con equipaggio che transitavano da sud.

Per una brezza, una corrente oceanica o un uccello migratore, l'Arco non
era altro che una coppia di pilastri immobili tra l'oceano Indiano e il
golfo del Bengala. Si muovevano liberamente attorno e dentro lo spazio
dell'Arco come ogni nave che viaggiava da nord a sud.

Ma attraversando l'equatore su una nave da sud a novanta gradi a est di
Greenwich, ci si ritrovava a guardare indietro verso l'Arco da un mare
sconosciuto sotto uno strano cielo, a una quantità imprecisata di anni
luce di distanza dalla Terra.

Nella città di Madras un servizio di crociere ambizioso se non del tutto
legale aveva prodotto una serie di poster in lingua inglese che
annunciavano: {VIAGGIO FACILE SU UN PIANETA AMICHEVOLE}! L'Interpol
aveva fatto chiudere l'azienda - in quel periodo le Nazioni Unite
stavano ancora cercando di regolamentare i passaggi - ma i poster ci
avevano azzeccato. Com'era possibile? Solo gli Ipotetici conoscevano la
risposta.

Il divorzio di Diane era ufficiale, così mi disse, ma non aveva né
lavoro né prospettive. «Pensavo di raggiungerti\ldots» Sembrava esitante
e per niente simile a un Quarto, o a ciò che immaginavo dovesse
somigliargli. «Se per te va bene. Sinceramente ho bisogno d'aiuto per
trovare un posto e, capisci, sistemarmi.»

Quindi le trovai un lavoro nella clinica e presentai i documenti per
l'immigrazione. Mi raggiunse a Montreal quell'autunno stesso.

ooo

Fu un corteggiamento sottile, lento, all'antica (o semimarziano, forse),
durante il quale Diane e io ci riscoprimmo in modi del tutto nuovi. Non
eravamo più sotto la camicia di forza dello Spin né eravamo bambini in
cerca di conforto. Ci innamorammo, finalmente, come adulti.

Furono gli anni in cui la popolazione globale raggiunse gli otto
miliardi. Gran parte di tale crescita era stata incanalata nelle
megalopoli in espansione: Shanghai, Jakarta, Manila, la Cina costiera;
Lagos, Kinshasa, Nairobi, Maputo; Caracas, La Paz, Tegucigalpa - tutte
le conigliere del mondo sigillate dal fuoco e dallo smog. Ci sarebbero
voluti una decina di Archi per ridurre la crescita della popolazione, ma
l'affollamento provocò un'ondata ininterrotta di emigranti e
``pionieri'', molti dei quali ammassati nei compartimenti di carico
d'imbarcazioni illegali e consegnati alla riva di Port Magellan in fin
di vita o già morti.

Port Magellan fu il primo insediamento ad avere un nome. Oggi gran parte
di quel mondo è stata grosso modo mappata, principalmente per via aerea.
Port Magellan era sulla punta orientale di un continente che alcuni
avevano preso a chiamare ``Equatoria''. C'era anche una seconda massa
continentale (``Borea''), ancora più grande, che ne percorreva il polo
nord e si estendeva alla zona temperata del pianeta. I mari del Sud
erano ricchi di isole e arcipelaghi.

Il clima era benigno, l'aria fresca, la gravità al 95,5\% di quella
terrestre. Entrambi i continenti erano granai ricolmi. I mari e i fiumi
brulicavano di pesci. La leggenda che circolava nei bassifondi di Douala
e Kabul narrava che potevi prendere da mangiare dagli alberi giganti di
Equatoria e rifugiarti per dormire tra le loro accoglienti radici.

Non potevi. Port Magellan era un'enclave delle Nazioni Unite sorvegliata
dai soldati. Le baraccopoli proliferate nei dintorni erano senza
controllo e pericolose. Ma i funzionali villaggi di pescatori segnavano
la linea costiera per centinaia di chilometri; c'erano alberghi
turistici in costruzione attorno alle lagune di Reach Bay e Aussie
Harbor; e la prospettiva di terre fertili gratuite aveva spinto i
colonizzatori verso l'entroterra, lungo le valli dei fiumi White e New
Irrawaddi.

Ma quell'anno la notizia più importante dal nuovo mondo fu la scoperta
del secondo Arco. Era ubicato a mezzo mondo di distanza dal primo,
vicino alle appendici meridionali della massa continentale boreale e
dietro vi era un mondo completamente nuovo - questo, secondo i primi
resoconti, era un po' meno invitante; o forse là era solo in corso la
stagione delle piogge.

ooo

«Devono esserci altre persone come me» disse Diane, cinque anni dopo
l'era dello Spin. «Mi piacerebbe incontrarle.»

Le avevo dato la copia degli archivi marziani, tradotti alla meglio e
salvati su una serie di memory card, e li aveva studiati attentamente
con la stessa intensità con cui aveva affrontato la poesia vittoriana e
gli scritti sul Nuovo Regno.

Se il lavoro di Jason aveva avuto successo, allora sì, c'erano
senz'altro altri Quarti sulla Terra. Ma rendere pubblica la loro
presenza sarebbe stato un biglietto di prima classe per un penitenziario
federale. Durante la crisi economica che aveva seguito la fine dello
Spin l'amministrazione Lomax aveva messo un freno di sicurezza nazionale
a tutto ciò che era marziano, e le agenzie di sicurezza domestiche di
Lomax godettero di poteri di polizia senza limiti.

«Non ci pensi mai?» chiese lei, timidamente.

A diventare anch'io un Quarto, intendeva. Iniettarmi nel braccio una
dose controllata di liquido trasparente da una delle fiale che tenevo in
una cassetta di sicurezza sul retro dell'armadio della nostra camera.
Certo che ci avevo pensato. Ci avrebbe reso ancora più simili.

Ma era ciò che volevo? Ero consapevole dello spazio invisibile, del
divario tra la sua condizione di Quarto e la mia umanità non modificata,
ma non mi faceva paura. Certe notti, guardando i suoi occhi solenni, mi
pareva addirittura un tesoro prezioso. Era il canyon che definiva il
ponte, e il ponte che avevamo costruito era gradevole e forte.

Mi carezzò la mano, le sue dita lisce sulla mia pelle ruvida, un sottile
promemoria del fatto che il tempo non si era mai fermato, che un giorno
avrei potuto necessitare del trattamento, anche se non lo avessi voluto.

«Non ancora» risposi.

«Quando?»

«Quando sarò pronto.»

ooo

Il Presidente Lomax fu rimpiazzato dal Presidente Hughes e poi dal
Presidente Chaykin, ma erano tutti veterani delle stesse politiche
dell'era dello Spin. Videro la biotecnologia marziana come la nuova
bomba atomica, almeno in potenza, e decisero che apparteneva solo a
loro, una minaccia brevettata. La prima spedizione diplomatica di Lomax
verso le Cinque Repubbliche era stata una richiesta di nascondere le
informazioni biotecnologiche dalle trasmissioni marziane non codificate
verso la Terra. Aveva giustificato la richiesta con argomentazioni
plausibili sull'effetto che tale tecnologia avrebbe potuto avere su un
mondo politicamente diviso e spesso violento - aveva citato la morte di
Wun Ngo Wen come esempio - e fino a quel momento i marziani erano stati
al gioco.

Ma persino questi contatti sterilizzati con Marte avevano seminato
discordia. Le politiche egualitarie delle Cinque Repubbliche avevano
fatto di Wun Ngo Wen una sorta di postumo emblema del nuovo movimento
globale laburista. (Era irritante vedere la faccia di Wun sulle
targhette portate dai lavoratori delle fabbriche d'abbigliamento nelle
zone industriali asiatiche o dagli operai che facevano i circuiti
integrati nelle \emph{maquiladoras} dell'America centrale - ma dubito
che gli sarebbe dispiaciuto.)

ooo

Diane passò il confine per andare al funerale di E.D. quasi undici anni
dopo il giorno in cui la salvai dal ranch Condon.

Apprendemmo la notizia della sua morte dai notiziari. Il necrologio
menzionava di sfuggita che l'ex moglie di E.D., Carol, era morta sei
mesi prima. Fu un triste colpo. Carol aveva smesso di rispondere alle
nostre chiamate quasi dieci anni prima. Troppo pericoloso, aveva detto.
Per lei era già abbastanza sapere che eravamo al sicuro. E non c'era
nulla, in realtà, da dire.

(Diane andò sulla tomba della madre quando si trovava a Washington. Ciò
che la rese maggiormente triste, confessò, fu che la vita di Carol fosse
stata così incompleta: un verbo senza complemento oggetto, una lettera
anonima fraintesa per mancanza di una firma. «Non mi manca tanto lei
quanto ciò che avrebbe potuto essere.»)

Al funerale di E.D., Diane fece in modo di non identificarsi. Erano
presenti troppi amici del governo, compreso il procuratore generale e il
Vicepresidente in carica. Ma la sua attenzione ricadde su una donna
anonima tra le panche, con cui scambiò, di nascosto dagli altri, alcuni
sguardi: «Sapevo che era un Quarto» mi disse Diane. «Non ti so dire
esattamente bene come. Il suo portamento, quella sorta di aspetto senza
età, ma ancor di più era come se un segnale andasse avanti e indietro
tra noi.» E quando la cerimonia fu terminata Diane si avvicinò alla
donna per chiederle come avesse conosciuto E.D.

«Non lo conoscevo,» rispose la donna, «non di persona. Feci un periodo
come ricercatrice alla Perielio, ai tempi di Jason Lawton. Mi chiamo
Sylvia Tucker.»

Quel nome mi fece suonare un campanello, quando Diane me lo ripeté.
Sylvia Tucker era una degli antropologi che avevano lavorato con Wun Ngo
Wen al complesso in Florida. Era stata più amichevole di tutti gli altri
accademici assunti ed era possibile che Jase si fosse fidato di lei.

«Ci siamo scambiate l'indirizzo di posta elettronica» mi informò Diane.
«Nessuna di noi due ha detto la parola ``Quarto''. Ma entrambe
\emph{sapevamo}. Ne sono certa.»

Non ne seguì alcuna corrispondenza, ma di tanto in tanto Diane riceveva
dall'indirizzo di Sylvia Tucker ritagli di giornale digitalizzati,
riguardanti, per esempio: ``\emph{Chimico industriale arrestato a Denver
	con ordinanza di sicurezza e trattenuto a tempo indeterminato}'' o
``\emph{Clinica geriatrica a Mexico City chiusa per ordine federale}'' o
``\emph{Professore di sociologia all'Università della California rimane
	ucciso in un incendio, sospetto incendio doloso}.''

E così via.

Mi ero premurato di non tenere una lista dei nomi e degli indirizzi a
cui Jason aveva spedito i suoi pacchetti finali, e non li avevo neppure
memorizzati. Ma alcuni dei nomi citati dagli articoli mi parevano
verosimilmente familiari.

«Ci sta dicendo che sono braccati» disse Diane. «Il governo sta dando la
caccia ai Quarti.»

Per un mese discutemmo sul da farsi nel caso in cui avessimo attirato lo
stesso tipo di attenzioni. Dato l'apparato di sicurezza globale ordito
da Lomax e coloro che lo avevano sostituito, dove saremmo scappati?

Ma in realtà c'era una sola risposta plausibile. Un unico luogo dove
l'apparato non riusciva a operare e la sorveglianza era del tutto cieca.
Perciò facemmo i nostri programmi - quei passaporti, quel conto
bancario, quella rotta attraverso l'Europa fino all'Asia meridionale - e
li mettemmo da parte finché non ce ne fosse stato bisogno.

Poi Diane ricevette un'ultima comunicazione da Sylvia Tucker, una
singola parola:

``Vai'' diceva.

E noi andammo.

ooo

Sull'ultimo volo del viaggio, arrivando a Sumatra, Diane disse: «Sei
sicuro di volerlo fare?»

Avevo preso la decisione giorni prima, durante uno scalo ad Amsterdam,
quando avevamo ancora paura che ci stessero seguendo, che i nostri
passaporti potessero essere stati segnalati o che il nostro
approvvigionamento di farmaci marziani venisse confiscato.

«Sì,» annuii, «ora. Prima di andare dall'altra parte.»

«Sei sicuro?»

«Non sarò mai più sicuro di così.»

No, non ero sicuro. Ma disposto. Disposto, alla fine, a cogliere ciò che
poteva essere perso, disposto ad abbracciare ciò che potevo ottenere.

E così affittammo una camera a Padang, al terzo piano di un hotel in
stile coloniale, dove per un po' saremmo passati inosservati.

``Tutti cadiamo,'' dissi tra me, ``e finiamo a terra da qualche parte.''
